\documentclass[11pt,a4paper]{article}
% Shared preamble for RuPhO English problem statements (match T1 2026 style).
% Use: \input{latex/rupho_en_preamble.tex} after \documentclass.
\usepackage[margin=1in]{geometry}
\usepackage{amsmath,amssymb}
\usepackage{graphicx}
\usepackage{float}
\usepackage{enumitem}
\usepackage{microtype}
\usepackage{caption}
\usepackage{tabularx}
\usepackage{hyperref}

\captionsetup{font=small,labelfont=bf,skip=6pt}

\setlist[itemize]{leftmargin=1.2cm}
\setlist[enumerate]{leftmargin=1.2cm}

% One outer box per subproblem: [id | statement | points] (no inner rules)
\newcommand{\problembox}[3]{%
  \par\addvspace{0.42em}%
  \noindent
  \setlength{\fboxsep}{7.5pt}%
  \setlength{\fboxrule}{0.95pt}%
  \fbox{%
    \begin{minipage}{\dimexpr\linewidth-2\fboxsep-2\fboxrule\relax}
    \begin{tabularx}{\linewidth}{@{}>{\raggedright\arraybackslash\bfseries}p{2.1em} @{\hspace{0.9em}} >{\raggedright\arraybackslash}X @{\hspace{0.9em}} >{\raggedleft\arraybackslash\bfseries}p{2.4em}@{}}
    #1 & #2 & #3
    \end{tabularx}
    \end{minipage}%
  }%
  \par\addvspace{0.32em}%
}

\title{\textbf{Proton in a capacitor}}
\author{\normalsize X16 (2016) --- English translation}
\date{}
\begin{document}
\maketitle

In the middle of a flat capacitor with a distance between the plates $d=2~\text{cm}$ there is a proton. Plate No. 1 of the capacitor is grounded, and at some point in time, alternating voltage $U(t)=U_0 \sin \omega t$, where $U_0=1~\text{mV}$, $\omega = 10^5~\text{s}^{-1}$, is applied to plate No. 2. After some time $\tau$ the proton collides with one of the plates.

\problembox{A1}{Which (1st or 2nd) plate did the proton collide with? Rate the time $\tau$.}{3.00}

Ignore the force of gravity and the field of induction charges. Proton charge $e = 1.6 \cdot 10^{-19}~\text{C}$, mass $m = 1.67 \cdot 10^{-27}~\text{kg}$.

\vspace{1em}
\noindent\textit{Source:} \url{https://pho.rs/p/3999}

\end{document}